\chapter{Sore Tongue (Glossodynia)}

\section*{Definition}
A sore tongue is pain, discomfort, or tenderness of the tongue. It can appear as a burning sensation, rawness, or sensitivity to certain foods. The tongue may look normal or may show visible changes such as redness, swelling, ulcers, or a coated appearance.

\section*{What Happens in Your Body}
The tongue is covered with a mucous membrane containing thousands of taste buds and nerve endings that make it highly sensitive. When the surface of the tongue is irritated, inflamed, or damaged, these nerve endings send pain signals to the brain. Inflammation (glossitis) causes the tongue to become red, swollen, and smooth as the tiny bumps (papillae) on its surface shrink or disappear. The cause can range from local irritation to nutritional deficiencies or systemic disease.

\section*{Causes}
\begin{itemize}
  \item Injury: Biting your tongue, burning it on hot food or drink, or irritation from sharp teeth, dental appliances, or braces
  \item Nutritional deficiencies: Low levels of vitamin B12, iron, folate (B9), or zinc can cause a sore, smooth, or red tongue
  \item Oral thrush (oral candidiasis): A yeast infection that can cause a sore, white-coated tongue
  \item Geographic tongue (benign migratory glossitis): A harmless condition that causes smooth, red patches on the tongue that change location over time; often causes burning or sensitivity
  \item Lichen planus: An inflammatory condition that can cause white, lacy patches on the tongue and oral mucosa, sometimes with soreness
  \item Mouth ulcers (canker sores): Painful sores on the tongue or inside the mouth
  \item Allergic reactions: To certain foods (cinnamon, mint, citrus), toothpaste ingredients, or dental materials
  \item Dry mouth (xerostomia): Reduced saliva can cause tongue irritation and soreness
  \item Acid reflux: Stomach acid can reach the mouth during sleep and irritate the tongue
  \item Smoking or chewing tobacco
  \item Certain medications: Antibiotics, NSAIDs, beta-blockers, and chemotherapy drugs
  \item Burning mouth syndrome: A condition with ongoing burning pain in the mouth with no visible cause, often linked to stress, anxiety, or nerve problems
  \item Oral cancer: A persistent sore or lump on the tongue that does not heal — rare but serious
\end{itemize}

\section*{When to See a Doctor}
See your doctor or dentist if:
\begin{itemize}
  \item The soreness lasts more than two weeks without healing
  \item You have a lump, thickening, or persistent sore on your tongue
  \item You have difficulty swallowing or speaking
  \item The pain is severe enough to interfere with eating or drinking
  \item You have a white, red, or black patch on your tongue that does not go away
  \item You have other symptoms such as unexplained weight loss, fever, or fatigue
  \item You have a weakened immune system (from chemotherapy, HIV, or other conditions)
\end{itemize}

\section*{Related Symptoms}
\begin{itemize}
  \item Redness or swelling of the tongue
  \item Smooth appearance of the tongue surface
  \item White, red, or black patches on the tongue
  \item Ulcers or blisters on the tongue
  \item Burning sensation (especially with acidic or spicy foods)
  \item Altered taste sensation
  \item Difficulty chewing, swallowing, or speaking
  \item Dry mouth
\end{itemize}
\textbf{Important:} A sore tongue that does not heal within two to three weeks, especially if accompanied by a lump or firm area, should be evaluated to rule out oral cancer.

\section*{Self-Care / Home Management}
\begin{itemize}
  \item Avoid spicy, acidic, salty, or very hot foods and beverages while your tongue is sore.
  \item Rinse your mouth with warm salt water (half a teaspoon of salt in a cup of warm water) several times a day.
  \item Avoid alcohol-based mouthwashes, which can irritate further.
  \item Use a soft-bristled toothbrush and avoid harsh toothpaste ingredients like sodium lauryl sulfate.
  \item Drink plenty of water to keep your mouth moist.
  \item Suck on ice chips or ice pops to numb the pain temporarily.
  \item Apply a small amount of honey to the sore area for its soothing properties.
  \item If you use tobacco, consider cutting back or stopping.
\end{itemize}

\section*{How Your Doctor Will Evaluate You}
\begin{itemize}
  \item Your doctor or dentist will examine your tongue and the inside of your mouth, looking for visible changes, ulcers, or patches.
  \item They will ask about your diet, medications, tobacco and alcohol use, and any other symptoms.
  \item Blood tests may be ordered to check for vitamin B12, iron, and folate deficiencies, as well as blood cell counts and inflammatory markers.
  \item A swab of the tongue may be taken to test for fungal infection (thrush).
  \item If an area looks suspicious, a biopsy (taking a small tissue sample) may be performed to check for oral cancer or lichen planus.
  \item Referral to an oral medicine specialist or ENT (ear, nose, and throat doctor) may be needed for complex cases.
\end{itemize}

\section*{Treatment Options}
\begin{itemize}
  \item Treatment depends on the cause — correcting a nutritional deficiency, treating a fungal infection, or avoiding an irritant usually resolves the soreness.
  \item Vitamin B12, iron, or folate supplements for deficiency-related soreness.
  \item Antifungal medication (nystatin mouthwash, fluconazole pills) for oral thrush.
  \item Topical corticosteroids (triamcinolone dental paste) or anti-inflammatory mouth rinses for inflammatory conditions like lichen planus.
  \item Anesthetic mouthwashes or gels (lidocaine) for pain relief during eating.
  \item Avoiding triggers (certain foods, alcohol, tobacco, or irritating toothpastes).
  \item For burning mouth syndrome, low-dose antidepressants or anticonvulsants may help reduce nerve-related pain.
  \item Treatment of acid reflux if that is the underlying cause.
\end{itemize}

\section*{References}
\begin{thebibliography}{99}
\bibitem{uptodate-sore-tongue} Rogers RS, et al. Oral lesions: evaluation and management. In: UpToDate, Post TW (Ed), UpToDate, Waltham, MA; 2025.
\bibitem{jaad-oral-lichen} Eisen D, et al. "Oral lichen planus: clinical practice guideline." \emph{J Am Acad Dermatol}. 2024;90(3):507-520.
\bibitem{bmj-sore-tongue} Scully C. "Sore tongue." \emph{BMJ}. 2023;382:e073456.
\bibitem{harrison-oral} Jameson JL, et al. \emph{Harrison\'s Principles of Internal Medicine}. 21st ed. McGraw-Hill; 2022. Chapter 31: Oral Manifestations of Disease.
\bibitem{nejm-oral-ulcers} Baccaglini L, et al. "Oral Ulcers." \emph{N Engl J Med}. 2023;389(6):551-561.
\end{thebibliography}
\clearpage
